\emph{First-order splitting.}
A twist by $\phi$ of the $U(1)$ symmetry generated by $m$ imposes $\phi\to\phi+2\pi$ boundary conditions on the operators of charge $m$ and shifts the momentum quantum number to $e+\phi/2\pi$ in Eq.~\eqref{eq:spectrum}~\cite{Alcaraz1988}.
At $\phi=\pi$ the lowest neutral states are $|{\pm\tfrac12}\rangle\coloneqq|e=\pm\tfrac12,m=0\rangle$, with dimension $K/4$, and $C$ exchanges them.
Since $V_{\pm1,0}$ changes $e$ by $\pm1$, the only nonzero matrix elements of $\Vd$ in this doublet are $\langle{\pm\tfrac12}|V_{\pm1,0}|{\mp\tfrac12}\rangle$, which equal the OPE coefficient of three vertex operators, $1$ in the normalization of Sec.~\ref{sec:model}.
In the basis $|C\!=\!\pm1\rangle=(|{+\tfrac12}\rangle\pm|{-\tfrac12}\rangle)/\sqrt2$, $\Vd=(V_{1,0}+V_{-1,0})/\sqrt2$ is therefore diagonal with eigenvalues $\pm1/\sqrt2$.
With the first-order shift $2\pi\gd$ times the matrix element on a cylinder of circumference $2\pi$ (Sec.~\ref{sec:k1}), the $C$-even level moves by $+\sqrt2\pi\gd$ and the $C$-odd level by $-\sqrt2\pi\gd$, which gives Eq.~\eqref{eq:split}.
The sign of $\gd$ entering Eq.~\eqref{eq:split} is tied to the phase convention of the vertex operators; as for $\kappa$ in Sec.~\ref{sec:compensation}, we fix it relative to the $a$/$b$ deformation, and both readouts use the same convention.
